% §10.4 MinHash and LSH
%
% Cross-reference resolution status:
%
% DEFINED here:
%   \label{sec:minhash-lsh}                       -- this section
%   \label{def:minhash}                           -- the construction
%   \label{alg:minhash-signature-estimate}        -- signature/estimate pseudocode
%   \label{prop:minhash-collision}                -- collision-probability lemma
%   \label{prop:minhash-error}                    -- Hoeffding bound proposition
%
% RESOLVED (chapter 9):
%   \Cref{ch:approximate-functions}               -- chapter 9
%   \Cref{def:n-ary-channel}                      -- §9.2
%   \Cref{thm:n-ary-kronecker}                    -- §9.3
%
% RESOLVED (chapter 10, prior sections):
%   \Cref{sec:sketches-as-bernoulli}              -- §10.1
%   \Cref{sec:countmin}                           -- §10.2
%   \Cref{sec:hyperloglog}                        -- §10.3
%   \Cref{prop:hyperloglog-error}                 -- §10.3
%
% FORWARD REFS within chapter 10 (resolve in Plan 10 Task 6):
%   \Cref{sec:ch10-bridge}                        -- §10.5
%
% BIBLIOGRAPHIC (resolved in §10.6, Plan 10 Task 7):
%   broder-1997                                   -- MinHash original
%   indyk-motwani-1998                            -- LSH original

\section{MinHash and LSH}
\label{sec:minhash-lsh}

\subsection*{Problem setup}

The similarity question on two sets is the third streaming primitive
of this chapter. Given two sets $A, B \subseteq \Universe$, the
\emph{Jaccard similarity} is
\[
  J(A, B) \;=\; \frac{|A \cap B|}{|A \cup B|},
\]
a value in $[0, 1]$ that is $0$ when the sets are disjoint and $1$
when they coincide. Exact computation of $J(A, B)$ requires a pass
through both sets and total work $O(|A| + |B|)$; storing the sets
themselves costs space linear in their sizes, which is a deal-breaker
when the inputs are documents over a vocabulary of millions of
shingles and the comparison must be made across a corpus of millions
of documents. The MinHash sketch (\textcite{broder-1997}) reduces
each set to a signature of $k$ small integers, with $k$ a tunable
accuracy parameter, and estimates $J(A, B)$ from the two signatures
alone in time $O(k)$. Pairwise similarity queries that would otherwise
be quadratic in the corpus size become quadratic in the signature
length, and the further extension to Locality-Sensitive Hashing (LSH)
reduces the corpus-wide comparison to a sub-linear approximate
nearest-neighbour search.

\subsection*{The MinHash trick}

Fix a hash $h : \Universe \to [0, 1]$ that is uniform across the
universe and injective on $A \cup B$ (no hash collisions among items
of interest). For a non-empty set $S \subseteq \Universe$, write
$\min h(S) = \min_{a \in S} h(a)$ for the smallest hash value attained
on $S$, and write $\arg\min h(S) = \arg\min_{a \in S} h(a)$ for the
item of $S$ at which that minimum is attained. The central
observation is the following.

\begin{proposition}[MinHash collision probability]
\label{prop:minhash-collision}
Let $A, B \subseteq \Universe$ be non-empty finite sets with
$A \cup B$ injectively hashed by $h$. Then
\[
  \Prob{\min h(A) = \min h(B)} \;=\; J(A, B).
\]
\end{proposition}

\begin{proof}
The minimum hash over $A \cup B$ is attained at a unique element
$a^* \in A \cup B$ (by injectivity). The event
$\{\min h(A) = \min h(B)\}$ holds exactly when $a^* \in A \cap B$:
if $a^*$ lies in both $A$ and $B$, then $a^*$ is the argmin of $h$
restricted to $A$ and also the argmin of $h$ restricted to $B$, and
the two minima coincide; if $a^*$ lies in only one of the two sets,
the other set's minimum is strictly larger. By the uniformity of $h$
across $A \cup B$, every element of $A \cup B$ is equally likely to
be the argmin, so
\[
  \Prob{a^* \in A \cap B}
  \;=\; \frac{|A \cap B|}{|A \cup B|}
  \;=\; J(A, B),
\]
which is the claimed identity.
\end{proof}

The proof is short because the identity is essentially combinatorial.
What the proof says, restated, is that the minimum hash over a union
is itself a uniform random sample from the union, and the probability
this sample comes from a particular subset is the subset's
proportion. A single hash $h$ converts the set-similarity question
into a Bernoulli coin flip with bias $J(A, B)$. Estimating $J$ is
then a matter of repeating the experiment.

\subsection*{The construction}

Pick $k$ hash functions $h_1, h_2, \ldots, h_k$, each independent and
each uniform across $\Universe$. For each set $A \subseteq \Universe$,
compute the \emph{MinHash signature}
\[
  \sigma(A)
  \;=\; \bigl(\sigma_1(A), \sigma_2(A), \ldots, \sigma_k(A)\bigr)
  \;\in\; [0, 1]^k,
\]
where $\sigma_i(A) = \min_{a \in A} h_i(a)$. The signature is a
length-$k$ vector that summarizes $A$: it discards the items of $A$
themselves and retains only the $k$ minima.

\begin{definition}[MinHash sketch]
\label{def:minhash}
The \emph{MinHash sketch} of parameter $k$ for a set $A \subseteq
\Universe$ is the signature $\sigma(A) = (\sigma_1(A), \ldots,
\sigma_k(A))$ where each $\sigma_i(A) = \min_{a \in A} h_i(a)$ for
independent uniform hashes $h_i$. The \emph{MinHash estimator} of the
Jaccard similarity $J(A, B)$ is
\[
  \hat J(A, B)
  \;=\; \frac{1}{k} \,
        \bigl|\{i \in \{1, \ldots, k\} : \sigma_i(A) = \sigma_i(B)\}\bigr|,
\]
the fraction of signature components on which the two sets agree.
\end{definition}

\Cref{alg:minhash-signature-estimate} states the signature
construction and the pairwise estimator in pseudocode.

\begin{algorithm}[ht]
  \caption{MinHash signature and Jaccard estimate}
  \label{alg:minhash-signature-estimate}
  \SetKwInOut{Input}{input}
  \SetKwInOut{Output}{output}
  \SetKwProg{func}{procedure}{}{}
  \SetKwProg{fn}{function}{}{}

  \fn{\textsc{Signature}($A$)}{
    \For{$i \gets 1$ \KwTo $k$}{
      $\sigma_i \gets \min_{a \in A} h_i(a)$\;
    }
    \Return $(\sigma_1, \sigma_2, \ldots, \sigma_k)$\;
  }
  \BlankLine
  \fn{\textsc{Estimate}($\sigma(A), \sigma(B)$)}{
    $c \gets 0$\;
    \For{$i \gets 1$ \KwTo $k$}{
      \lIf{$\sigma_i(A) = \sigma_i(B)$}{$c \gets c + 1$}
    }
    \Return $c / k$\;
  }
\end{algorithm}

The signature for $A$ takes $O(k \cdot |A|)$ time to compute (one
pass per hash function) and $O(k)$ space to store. The pairwise
estimator takes $O(k)$ time, independent of the original set sizes.
The signature is also \emph{streaming}: items can be inserted into
$A$ one at a time, with each insertion updating each $\sigma_i$ in
$O(1)$ via $\sigma_i \gets \min(\sigma_i, h_i(a))$. Deletion is not
supported, which is the price of the min-based summary; an item that
once attained the minimum cannot be removed without rescanning the
set.

\subsection*{Error bound}

Each component of the signature gives a Bernoulli trial for the
Jaccard event. Let $Y_i = \mathbf{1}[\sigma_i(A) = \sigma_i(B)]$ for
$i = 1, \ldots, k$. By \Cref{prop:minhash-collision}, each $Y_i$ is
Bernoulli with parameter $J = J(A, B)$, and the $Y_i$ are independent
because the hash functions $h_i$ are independent. The estimator is
\[
  \hat J \;=\; \frac{1}{k} \sum_{i=1}^k Y_i,
\]
the sample mean of $k$ independent Bernoulli($J$) variables.

\begin{proposition}[MinHash error bound]
\label{prop:minhash-error}
For a MinHash sketch of parameter $k$ built from independent uniform
hashes, the estimator $\hat J$ of the Jaccard similarity $J = J(A, B)$
satisfies
\[
  \mathbb{E}[\hat J] \;=\; J,
  \qquad
  \mathrm{Var}[\hat J] \;=\; \frac{J(1 - J)}{k},
\]
and, by Hoeffding's inequality, for every $\epsilon > 0$,
\[
  \Prob{|\hat J - J| \;\ge\; \epsilon}
  \;\le\; 2 \exp\bigl(-2 k \epsilon^2\bigr).
\]
\end{proposition}

\begin{proof}
Linearity of expectation gives $\mathbb{E}[\hat J] = \mathbb{E}[Y_1]
= J$, and independence of the $Y_i$ gives $\mathrm{Var}[\hat J] =
\mathrm{Var}[Y_1] / k = J(1 - J)/k$. Hoeffding's inequality for a
sum of $k$ independent random variables each bounded in $[0, 1]$
states that the deviation of the sample mean from its expectation
exceeds $\epsilon$ in absolute value with probability at most
$2 \exp(-2 k \epsilon^2)$, which is the claimed tail bound applied to
$\hat J$.
\end{proof}

The bound is symmetric in over-estimation and under-estimation: the
MinHash estimator is two-sided in the same sense that the
HyperLogLog estimator is two-sided
(\Cref{prop:hyperloglog-error}), and unlike CountMin, which is
one-sided. Inverting the Hoeffding bound for a target failure
probability $\delta$ gives an additive accuracy $\epsilon = \sqrt{\ln
(2/\delta) / (2 k)}$, so $\hat J$ falls within $\pm \epsilon$ of $J$
with probability at least $1 - \delta$ when $k \ge \ln(2/\delta) /
(2 \epsilon^2)$. The signature length grows linearly in $\ln(1/\delta)$
and quadratically in $1/\epsilon$.

\subsection*{The $\Bernoulli{T}$ type}

MinHash instantiates the framework of
\Cref{ch:approximate-functions} with output type $T = [0, 1]$, the
unit interval of similarity values. The query is a function from
pairs of signatures to similarity estimates, and its error behaviour
is summarized by a channel matrix $Q$ in the sense of
\Cref{def:n-ary-channel}: rows are indexed by the true Jaccard
$J \in [0, 1]$ and columns by the reported estimate
$\hat J \in \{0, 1/k, 2/k, \ldots, 1\}$, with entry
\[
  Q_{J, \hat J}
  \;=\; \Prob{\textsc{Estimate}(\sigma(A), \sigma(B)) = \hat J
              \Given J(A, B) = J}.
\]
Because $k \hat J$ is the sum of $k$ independent Bernoulli($J$)
variables, each row of $Q$ is a Binomial$(k, J)$ distribution on the
$k+1$ possible values of $\hat J$:
\[
  Q_{J, j/k}
  \;=\; \binom{k}{j} J^j (1 - J)^{k - j},
  \qquad j = 0, 1, \ldots, k.
\]
The matrix is row-stochastic by construction and is symmetric in the
following sense: the row conditional on true similarity $J$ is the
mirror image, under the substitution $\hat J \mapsto 1 - \hat J$, of
the row conditional on true similarity $1 - J$. The contrast with
CountMin's upper-triangular $Q$ and with HyperLogLog's approximately
log-normal rows is that MinHash's $Q$ has a closed binomial form. The
two-sided concentration of \Cref{prop:minhash-error} is the
chapter~9 channel-matrix view of the Binomial distribution's
concentration around its mean.

\subsection*{LSH for approximate nearest neighbour}

Estimating $J(A, B)$ for a single pair $(A, B)$ from signatures of
length $k$ takes $O(k)$ time. Doing the same for every pair in a
corpus of $N$ documents takes $O(k N^2)$ time, which is intractable
once $N$ is in the millions. Locality-Sensitive Hashing
(\textcite{indyk-motwani-1998}) is a structural device that uses the
MinHash signatures themselves as bucket keys, so that pairs with high
Jaccard similarity collide into the same bucket with high
probability and pairs with low similarity do not. The corpus-wide
similarity search reduces to a bucket lookup: for a query document
$Q$, hash its signature into the same scheme and inspect only the
documents sharing a bucket with $Q$.

The construction partitions the length-$k$ MinHash signature into
$b$ \emph{bands} of $r$ \emph{rows} each, with $k = b r$. For each
band, the $r$ signature entries are concatenated into a single key
and hashed into a bucket. Two sets $A$ and $B$ become \emph{candidate
neighbours} if they share a bucket in at least one of the $b$ bands.
The probability that a single band agrees on all $r$ rows is $J^r$
(each row agrees independently with probability $J$ by
\Cref{prop:minhash-collision}), so the probability that at least one
band agrees is
\[
  p(J) \;=\; 1 - (1 - J^r)^b.
\]
The shape of $p$ as a function of $J$ is the celebrated S-curve of
LSH: small $J$ gives $p(J) \approx b J^r$, which is close to zero
when $J^r$ is small; large $J$ gives $p(J)$ close to one. The
transition is sharp, with the steep portion concentrated around the
threshold $J^* \approx (1/b)^{1/r}$. Choosing $(b, r)$ tunes the
threshold: a larger $r$ steepens the curve and pushes the threshold
right (more selective), and a larger $b$ raises the entire curve
(more permissive).

For approximate nearest neighbour with target similarity $J^*$, the
operational recipe is to pick $(b, r)$ so that $p(J)$ is small below
$J^*$ and close to one above $J^*$, then build the $b$ bucket tables
once over the corpus. A query hashes its signature into each of the
$b$ band tables in $O(b)$ time and returns the union of the matching
buckets as the candidate set; an exact pairwise Jaccard (or a MinHash
estimate) on the small candidate set then filters down to the true
neighbours. The expected candidate-set size is governed by the
corpus's similarity distribution and the chosen $(b, r)$, and is
typically a small constant times the number of true neighbours when
the parameters are tuned to the target $J^*$. The resulting query
time is sub-linear in $N$: the bucket lookup is $O(b)$, and the
candidate-filtering work is governed by the actual neighbour density,
not the corpus size.

\subsection*{A worked example}

Take two sets $A$ and $B$ with $|A| = |B| = 100$ and overlap
$|A \cap B| = 30$. The union has $|A \cup B| = |A| + |B| - |A \cap B|
= 170$, and the true Jaccard similarity is
\[
  J \;=\; \frac{30}{170} \;\approx\; 0.176.
\]
Build MinHash signatures of length $k = 200$ from independent uniform
hashes. By \Cref{prop:minhash-error}, the estimator $\hat J$ has
expectation $J \approx 0.176$ and variance
\[
  \mathrm{Var}[\hat J]
  \;=\; \frac{J(1 - J)}{k}
  \;=\; \frac{0.176 \cdot 0.824}{200}
  \;\approx\; 7.26 \cdot 10^{-4},
\]
so the standard error of $\hat J$ is approximately $0.027$. A single
sketch typically returns $\hat J$ within $\pm 0.027$ of $0.176$, and
the Hoeffding bound of \Cref{prop:minhash-error} gives, for
$\epsilon = 0.05$,
\[
  \Prob{|\hat J - J| \ge 0.05}
  \;\le\; 2 \exp(-2 \cdot 200 \cdot 0.0025)
  \;=\; 2 e^{-1}
  \;\approx\; 0.736,
\]
which is a loose bound at this scale because Hoeffding ignores the
$J(1 - J)$ factor that pulls the variance well below $1/4$ when $J$
is small. The standard-error estimate above, derived from the exact
binomial variance, is the practically relevant accuracy figure. A
representative realization (verified in the chapter notebook with
seed 0) returns $\hat J = 0.185$, a deviation of $0.009$ from the
true value, well inside one standard error. To halve the standard
error to $0.013$, quadruple the signature length to $k = 800$; this
is the same $1/\sqrt{k}$ space-accuracy scaling as HyperLogLog and
the same logarithmic-in-failure-probability dependence as CountMin,
adapted to the symmetric two-sided Bernoulli setting.

\subsection*{Looking ahead}

The MinHash construction closes the catalogue of three sketches that
this chapter builds from scratch. The next section,
\Cref{sec:ch10-bridge}, recaps the common pattern (construction,
$\Bernoulli{T}$ type, error bound, composition) across CountMin,
HyperLogLog, and MinHash, and hands off to chapter~13's type algebra,
where the composition rules informally invoked throughout this
chapter are axiomatized and proved in their full generality.
